This section shows how the instruction set maps onto common C library kernels. The examples are schematic assembly
sequences rather than ABI-mandated entry points.

The examples use R0 as a source pointer, R1 as a destination pointer, R2 as the REP counter, and R3 as a temporary
or fill-value register.

\subsection{REP Counter Convention}
REP and REPcc interpret the selected Rn register as an ordinary signed two's-complement counter. A zero counter
performs no iteration. Positive counters decrement toward zero; negative counters increment toward zero.

The same architectural Rn value is visible to indexed effective-address calculation. Address update is part of the
EA syntax; after each successful element access, the selected base register advances by the element size.

\subsection{memcpy}
The forward byte-copy form uses source and destination base pointers and a positive signed count.

@MEMCPY_EXAMPLE@

\subsection{memmove}
memmove chooses direction according to overlap. Non-overlapping and forward copies use the same shape as memcpy;
overlapping reverse copies use the matching predecrement address form.

@MEMMOVE_EXAMPLE@

\subsection{memset}
memset uses a register source and a memory destination. The counter and auto-update convention is the same as the
forward copy traversal.

@MEMSET_EXAMPLE@

\subsection{memcmp}
memcmp can use REPEQ to continue while the observed comparison result is equal. The final FLAGS value is the result
of the last completed comparison.

@MEMCMP_EXAMPLE@

\subsection{Bulk Character Extension}
The orthogonal EXTS/EXTZ memory forms let a C library or codec routine widen elements while copying.

@BULK_EXTENSION_EXAMPLE@
